% Options for packages loaded elsewhere
\PassOptionsToPackage{unicode}{hyperref}
\PassOptionsToPackage{hyphens}{url}
\PassOptionsToPackage{dvipsnames,svgnames,x11names}{xcolor}
%
\documentclass[
  11pt,
]{article}
\usepackage{amsmath,amssymb}
\usepackage{iftex}
\ifPDFTeX
  \usepackage[T1]{fontenc}
  \usepackage[utf8]{inputenc}
  \usepackage{textcomp} % provide euro and other symbols
\else % if luatex or xetex
  \usepackage{unicode-math} % this also loads fontspec
  \defaultfontfeatures{Scale=MatchLowercase}
  \defaultfontfeatures[\rmfamily]{Ligatures=TeX,Scale=1}
\fi
\usepackage{lmodern}
\ifPDFTeX\else
  % xetex/luatex font selection
\fi
% Use upquote if available, for straight quotes in verbatim environments
\IfFileExists{upquote.sty}{\usepackage{upquote}}{}
\IfFileExists{microtype.sty}{% use microtype if available
  \usepackage[]{microtype}
  \UseMicrotypeSet[protrusion]{basicmath} % disable protrusion for tt fonts
}{}
\makeatletter
\@ifundefined{KOMAClassName}{% if non-KOMA class
  \IfFileExists{parskip.sty}{%
    \usepackage{parskip}
  }{% else
    \setlength{\parindent}{0pt}
    \setlength{\parskip}{6pt plus 2pt minus 1pt}}
}{% if KOMA class
  \KOMAoptions{parskip=half}}
\makeatother
\usepackage{xcolor}
\usepackage[margin=0.9in]{geometry}
\usepackage{color}
\usepackage{fancyvrb}
\newcommand{\VerbBar}{|}
\newcommand{\VERB}{\Verb[commandchars=\\\{\}]}
\DefineVerbatimEnvironment{Highlighting}{Verbatim}{commandchars=\\\{\}}
% Add ',fontsize=\small' for more characters per line
\newenvironment{Shaded}{}{}
\newcommand{\AlertTok}[1]{\textcolor[rgb]{1.00,0.00,0.00}{\textbf{#1}}}
\newcommand{\AnnotationTok}[1]{\textcolor[rgb]{0.38,0.63,0.69}{\textbf{\textit{#1}}}}
\newcommand{\AttributeTok}[1]{\textcolor[rgb]{0.49,0.56,0.16}{#1}}
\newcommand{\BaseNTok}[1]{\textcolor[rgb]{0.25,0.63,0.44}{#1}}
\newcommand{\BuiltInTok}[1]{\textcolor[rgb]{0.00,0.50,0.00}{#1}}
\newcommand{\CharTok}[1]{\textcolor[rgb]{0.25,0.44,0.63}{#1}}
\newcommand{\CommentTok}[1]{\textcolor[rgb]{0.38,0.63,0.69}{\textit{#1}}}
\newcommand{\CommentVarTok}[1]{\textcolor[rgb]{0.38,0.63,0.69}{\textbf{\textit{#1}}}}
\newcommand{\ConstantTok}[1]{\textcolor[rgb]{0.53,0.00,0.00}{#1}}
\newcommand{\ControlFlowTok}[1]{\textcolor[rgb]{0.00,0.44,0.13}{\textbf{#1}}}
\newcommand{\DataTypeTok}[1]{\textcolor[rgb]{0.56,0.13,0.00}{#1}}
\newcommand{\DecValTok}[1]{\textcolor[rgb]{0.25,0.63,0.44}{#1}}
\newcommand{\DocumentationTok}[1]{\textcolor[rgb]{0.73,0.13,0.13}{\textit{#1}}}
\newcommand{\ErrorTok}[1]{\textcolor[rgb]{1.00,0.00,0.00}{\textbf{#1}}}
\newcommand{\ExtensionTok}[1]{#1}
\newcommand{\FloatTok}[1]{\textcolor[rgb]{0.25,0.63,0.44}{#1}}
\newcommand{\FunctionTok}[1]{\textcolor[rgb]{0.02,0.16,0.49}{#1}}
\newcommand{\ImportTok}[1]{\textcolor[rgb]{0.00,0.50,0.00}{\textbf{#1}}}
\newcommand{\InformationTok}[1]{\textcolor[rgb]{0.38,0.63,0.69}{\textbf{\textit{#1}}}}
\newcommand{\KeywordTok}[1]{\textcolor[rgb]{0.00,0.44,0.13}{\textbf{#1}}}
\newcommand{\NormalTok}[1]{#1}
\newcommand{\OperatorTok}[1]{\textcolor[rgb]{0.40,0.40,0.40}{#1}}
\newcommand{\OtherTok}[1]{\textcolor[rgb]{0.00,0.44,0.13}{#1}}
\newcommand{\PreprocessorTok}[1]{\textcolor[rgb]{0.74,0.48,0.00}{#1}}
\newcommand{\RegionMarkerTok}[1]{#1}
\newcommand{\SpecialCharTok}[1]{\textcolor[rgb]{0.25,0.44,0.63}{#1}}
\newcommand{\SpecialStringTok}[1]{\textcolor[rgb]{0.73,0.40,0.53}{#1}}
\newcommand{\StringTok}[1]{\textcolor[rgb]{0.25,0.44,0.63}{#1}}
\newcommand{\VariableTok}[1]{\textcolor[rgb]{0.10,0.09,0.49}{#1}}
\newcommand{\VerbatimStringTok}[1]{\textcolor[rgb]{0.25,0.44,0.63}{#1}}
\newcommand{\WarningTok}[1]{\textcolor[rgb]{0.38,0.63,0.69}{\textbf{\textit{#1}}}}
\usepackage{longtable,booktabs,array}
\usepackage{calc} % for calculating minipage widths
% Correct order of tables after \paragraph or \subparagraph
\usepackage{etoolbox}
\makeatletter
\patchcmd\longtable{\par}{\if@noskipsec\mbox{}\fi\par}{}{}
\makeatother
% Allow footnotes in longtable head/foot
\IfFileExists{footnotehyper.sty}{\usepackage{footnotehyper}}{\usepackage{footnote}}
\makesavenoteenv{longtable}
\setlength{\emergencystretch}{3em} % prevent overfull lines
\providecommand{\tightlist}{%
  \setlength{\itemsep}{0pt}\setlength{\parskip}{0pt}}
\setcounter{secnumdepth}{-\maxdimen} % remove section numbering
\usepackage{needspace}
\ifLuaTeX
  \usepackage{selnolig}  % disable illegal ligatures
\fi
\usepackage{bookmark}
\IfFileExists{xurl.sty}{\usepackage{xurl}}{} % add URL line breaks if available
\urlstyle{same}
\hypersetup{
  pdftitle={Completing the residual-and-repair construction for benzel tilings},
  colorlinks=true,
  linkcolor={Maroon},
  filecolor={Maroon},
  citecolor={Blue},
  urlcolor={Blue},
  pdfcreator={LaTeX via pandoc}}

\title{Completing the residual-and-repair construction for benzel
tilings}
\usepackage{etoolbox}
\makeatletter
\providecommand{\subtitle}[1]{% add subtitle to \maketitle
  \apptocmd{\@title}{\par {\large #1 \par}}{}{}
}
\makeatother
\subtitle{All deletion residues, the exceptional endpoints, and the
original Split-Zero comparison}
\author{}
\date{18 September 2026}

\begin{document}
\maketitle

{
\hypersetup{linkcolor=}
\setcounter{tocdepth}{3}
\tableofcontents
}
\section{1. Result, scope, and
provenance}\label{result-scope-and-provenance}

\textbf{Constructive residue-zero theorem.} For every pair of integers
\[
 d\ge2,\qquad 0\le h\le t_d,\qquad t_d=d(d-1)/2,
\] the constructor in \texttt{complete\_theorem.py} produces an original
type-103 tiling of \[
 W_h(d)=V(d+3h,2d+3h).
\] It contains exactly \[
 \Delta=t_d-h\quad\text{right stones},\qquad 3h(h+d)\quad\text{bones}.
\] Reflection produces the reversed parameter pair. Every coefficient in
the output is a nonnegative integer; the original incidence vector is
the all-ones vector of the original region. No search is called by the
constructor. All canonical deletion counts are covered, including 7, 8,
10, and 11.

This is a complete, certificate-assisted research proof of this
construction, not a declaration of independent specialist acceptance.
Its prerequisites are the retained full stable-packing argument, the
retained threefold construction, and the retained first nonzero residue
construction. Their original source files are included unchanged and
their exact verifiers are replayed. The present note supplies the
previously missing second residue construction and the four exceptional
endpoints, and proves that the parameter coverage is exhaustive.

The original Problem 7 has two other residue families. Section 11 gives
a direct construction for residue one and records the exact previously
established compression-theorem input for residue two. Thus the hard
residue-zero gap in the previously recorded route is closed. This
package does not pretend to supply a new self-contained proof of that
published compression theorem.

\textbf{Provenance correction.} The owner supplied a report of Akshat
Sharma's earlier full-solution claim, \emph{A Constructive Proof of
Propp's Problem 7}, dated 8 September 2026, and the earlier published
bone-homology calculation of Muzika Dizdarević, Timotijević, and
Živaljević. Neither bare tiling existence nor the bone quotient is
advertised here as new. This continuation uses the existing residual
packing and original local repairs; it did not import Sharma's
construction. It establishes our displayed argument, not a line-by-line
audit of Sharma's manuscript. Independent review and a new Lean build
have not been performed.

\Needspace{20\baselineskip}

\section{2. Original cells and inverse parameter
maps}\label{original-cells-and-inverse-parameter-maps}

The axial cell is \((x,y)\in\mathbb Z^2\), corresponding to the
barycentric cell \((x,y,1-x-y)\); projection onto the first two
coordinates is the inverse. Its differences are \[
 u=y-x,\qquad v=1-x-2y,\qquad w=2x+y-1.
\] The original benzel consists of cells with all three differences in
\([1-a,b-1]\). The four original permitted tiles are

\begin{longtable}[]{@{}ll@{}}
\toprule\noalign{}
Kind & Cell offsets at its original anchor \\
\midrule\noalign{}
\endhead
\bottomrule\noalign{}
\endlastfoot
Right stone \(R\) & \((0,0),(1,0),(0,1)\) \\
Horizontal bone \(H\) & \((0,0),(1,0),(2,0)\) \\
Vertical bone \(V\) & \((0,0),(0,1),(0,2)\) \\
Diagonal bone \(D\) & \((0,0),(1,-1),(2,-2)\) \\
\end{longtable}

The differential \(\partial\) sends each placed tile generator to the
sum of its three original cell generators. The same map on free
nonnegative monoids sends actual tilings to their original coverage
vectors. The inclusion of these monoids in their free integer group
completions commutes with \(\partial\); no integer chain is called a
tiling merely because it has the same image.

Use the original affine rotation and reflection \[
 \rho(x,y)=(y,1-x-y),\qquad F(x,y)=(x,1-x-y).
\] Their inverses are \(\rho^2\) and \(F\). On the original differences,
\[
 (u,v,w)\circ\rho=(v,w,u),\qquad
 (u,v,w)\circ F=(-w,-v,-u).
\] Consequently \(\rho\) preserves the original benzel, and \(F\) maps
\(V(a,b)\) bijectively to \(V(b,a)\). It preserves the right stone and
permutes the three bones. Each induced tile map and cell map commutes
with incidence by summing the same three cell images.

For the residue-zero family the original parameter maps are \[
 (d,h)\longmapsto(a,b)=(d+3h,2d+3h),\qquad
 d=b-a,\quad h=(2a-b)/3.
\] These are inverse on the stated integer lane. The right-minus-left
stone invariant becomes \(\Delta=t_d-h\). The region contains \[
 |W_h(d)|=3t_d+(9d-3)h+9h^2
          =3\Delta+9h(h+d)
\] original cells. The retained packing proof derives this count on the
original carrier; the finite verifier also reconstructs regions through
the integer inverse \[
 x=(1-v-2u)/3,\qquad y=x+u,
\] retaining its divisibility condition.

\section{3. The canonical finite core and its receiving
maps}\label{the-canonical-finite-core-and-its-receiving-maps}

For \(\Delta>0\), let \(m\) be the unique integer with \[
 t_{m-1}<\Delta\le t_m,
\] and set \(q=t_m-\Delta\). Then \[
 0\le q\le m-2,\qquad m\le d,\qquad
 h=t_d-t_m+q.
\] Conversely, these equations recover the original \(\Delta\) and
\(h\). Uniqueness follows from strict increase of the triangular numbers
from \(t_1=0\). For \(\Delta=0\), use \(m=1,q=0\). These choices are
implemented by an integer square-root inverse, with its defining
inequalities checked.

The full stable-packing proof is retained at
\texttt{../turn8/sources/turn6-PROOF.md}. Its original rank
representatives are \[
 p_{r,s}=(r-s,-s),\quad 0\le s<r,\quad
 \operatorname{rank}(p_{r,s})=t_r+s+1,
\] with their three rotated copies. Locating the unique minimum among
the three differences and then its negative value \(r\) and coordinate
\(s\) gives the inverse cell map. The first \(\Delta\) rank orbits form
the actual residual \(\Omega_\Delta\).

The retained positive band family \(\Gamma_\Delta(d,h)\) obeys \[
 \partial\Gamma_\Delta(d,h)+\mathbf1_{\Omega_\Delta}
 =\mathbf1_{W_h(d)}.                                      \tag{1}
\] Every band is made of original bones, and distinct bands and sectors
are disjoint. Its exact coordinate proof is included, rather than only a
reference to an unavailable construction. For triangular \(\Delta=t_m\),
\[
 \Omega_{t_m}=W_0(m),
\] and the explicit all-stone base is \[
 T_0(m)=\{R(m-2-2r-s,r-s):r,s\ge0,\ r+s\le m-2\}.          \tag{2}
\] For \(m=1\), this is empty. Equations (1) and (2) prove the entire
\(q=0\) family.

For \(q\ge1\), retain the common original source bones \[
 a_{p,y,\ell}=F\rho^pH(y-r_y-2-3\ell,y),
\quad
 r_y=\begin{cases}m-1&y\le q,\\m&y>q,\end{cases}            \tag{3}
\] where \[
 p\in\{0,1,2\},\quad1\le y\le m+q,\quad
 0\le\ell<\min(y,q).
\] Call this complete matching \(A\). It is disjoint from
\(\Omega_\Delta\).

At \(d=m\), \(h=q\). A prefix bone has index \(j=y-1-\ell\). In a tail
row, the original first anchor is \(y-3q-m+1\), and index \(j=q-1-\ell\)
gives exactly \(y-m-2-3\ell\). The index maps are \[
 j_{\rm tail}=j_{\rm stable}+q-y,
 \qquad j_{\rm stable}=j_{\rm tail}+y-q.                    \tag{4}
\] All original index bounds are retained. For \(d\ge m+1\),
\(h\ge m+q\), so the same source bones all belong to the finite prefix.
Equations (3) and (4) preserve the sector, anchor, orientation, and
three cells; hence they commute with the original differential.

A positive core tiling of \(\Omega_\Delta\sqcup\operatorname{cells}(A)\)
extends to the original benzel by adjoining the actual frozen bones
\(\Gamma_\Delta\setminus A\). Deleting those frozen bones is the inverse
on tilings containing them. This is the positive-fibre map used below;
the smallest exterior parameter has not been omitted.

\section{4. The uniform second-residue
scaffold}\label{the-uniform-second-residue-scaffold}

In this section \[
 k\ge4,\qquad m\ge3k+4,\qquad q=3k+2.
\] The same scaffold also admits the fixed specializations \(k=2,3\) in
Section 7.

Let \(B_8(m,k)\) and \(S_8(m,k)\) be the exact replacement bones and
retained base stones of the completed \(q=3k\) construction, whose full
sources and symbolic verifier are preserved in \texttt{../turn8}. Remove
from \(S_8\) the two original stones \[
 R(m-2-k,-k),\qquad \rho^2R(m-2-k,-k),                     \tag{5}
\] and translate the surviving stones and \(B_8\) by \((-1,-1)\). Denote
the resulting families by \(S\) and \(B_0\).

The two stones in (5) have base indices \((r,s)=(0,k)\) and
\((m-k-2,0)\). The domain gives \(m-k-2\ge2k+2\); they are distinct and
survive the earlier corner deletions. Their removal leaves exactly
\(t_m-3k-2\) stones. This choice is essential: deleting the first and
the once-rotated stone instead leaves the vacancy in the wrong place for
the second channel.

Define \(A_0\subset A\) by the following original band-label table. All
endpoints are included. An empty depth interval contributes no
generators.

\begin{longtable}[]{@{}
  >{\raggedright\arraybackslash}p{(\columnwidth - 4\tabcolsep) * \real{0.3333}}
  >{\raggedright\arraybackslash}p{(\columnwidth - 4\tabcolsep) * \real{0.3333}}
  >{\raggedright\arraybackslash}p{(\columnwidth - 4\tabcolsep) * \real{0.3333}}@{}}
\toprule\noalign{}
\begin{minipage}[b]{\linewidth}\raggedright
Sector
\end{minipage} & \begin{minipage}[b]{\linewidth}\raggedright
Inclusive depth interval
\end{minipage} & \begin{minipage}[b]{\linewidth}\raggedright
Inclusive row interval
\end{minipage} \\
\midrule\noalign{}
\endhead
\bottomrule\noalign{}
\endlastfoot
0 & \(0\) & \(1\le y\le m + 2k\) \\
0 & \(1\le\ell\le k - 1\) & \(\ell + 1\le y\le m + 2k - 2\) \\
0 & \(1\le\ell\le k - 1\) &
\(m + 2k + \ell - 1\le y\le m + 2k + \ell\) \\
0 & \(k\) & \(k + 1\le y\le m + 3k\) \\
0 & \(k + 1\le\ell\le 2k - 1\) & \(\ell + 1\le y\le 2k + 2\) \\
0 & \(2k\) & \(y=2k + 1\) \\
1 & \(0\le\ell\le k - 2\) & \(\ell + 1\le y\le m + 2k - 5\) \\
1 & \(0\le\ell\le k - 2\) &
\(m + 2k + \ell - 3\le y\le m + 2k + \ell - 2\) \\
1 & \(k - 1\) & \(k\le y\le m + 3k - 3\) \\
1 & \(k\) & \(k + 1\le y\le 2k - 1\) \\
1 & \(k\) & \(3k\le y\le 3k + 2\) \\
1 & \(k + 1\le\ell\le 2k - 2\) & \(\ell + 1\le y\le 2k - 1\) \\
2 & \(0\le\ell\le k - 3\) & \(\ell + 1\le y\le m + 2k - 5\) \\
2 & \(0\le\ell\le k - 3\) &
\(m + 2k + \ell - 2\le y\le m + 2k + \ell - 1\) \\
2 & \(k - 2\) & \(k - 1\le y\le m + 3k - 3\) \\
2 & \(k - 1\) & \(k\le y\le 2k - 1\) \\
2 & \(k - 1\) & \(3k\le y\le 3k + 2\) \\
2 & \(k\le\ell\le 2k - 3\) & \(\ell + 1\le y\le 2k - 1\) \\
\end{longtable}

The table is literal, with no chosen quotient coordinates replacing
\((p,y,\ell)\). Within each sector its overlapping depth ranges have
disjoint row intervals. All labels obey (3). \texttt{source\_domains.py}
reconstructs and exactly verifies the nonnegative lengths, membership
bounds, branch bounds, and pairwise distinctness over the complete
parameter domains.

Write the eighteen cells of the preceding first-residue scaffold as the
disjoint union \[
 Z_9=A_9\sqcup F_9\sqcup R_9\sqcup U_9,
\] where \[
 A_9=\alpha(2k-2,4-4k-m),\qquad
 \alpha(x,y)=\{(x,y),(x,y+1),(x+1,y-1)\},
\] \[
\begin{split}
 F_9=\{&(3k-1,1-m),(3k-1,2-m),(3k,2-m),\\
       &(3k,3-m),(3k+1,3-m),(3k+1,4-m)\},
\end{split}
\] \[
 R_9=R(m-k-4,1-k)
\] as an original three-cell set, and \[
 U_9=V(m+2k-5,2k-4)
 \sqcup\{(m+2k-2,2k-1),(m+2k-2,2k),(m+2k-1,2k)\}.
\] Here \(U_9\) has six cells. Define the actual affine maps \[
 T_1(x,y)=(x+1,y-2),\qquad T_2(x,y)=\rho^2(x,y)+(-2,1).
\] Their inverses are translation by \((-1,2)\) and
\(z\mapsto\rho(z-(-2,1))\). Their action on placed tiles is obtained by
applying the map to the three original cells.

The new 27-cell residual is \[
 Z_2=T_1(A_9\sqcup F_9\sqcup R_9)
       \sqcup T_2(F_9\sqcup R_9\sqcup U_9).                \tag{6}
\] The exact scaffold identity is \[
 \partial(B_0+S)+\mathbf1_{Z_2}
    =\mathbf1_{\Omega_\Delta}+\partial A_0.                \tag{7}
\] \texttt{source10.py} expresses (7) in six independent Laurent
variables. Its original integer numerator vanishes identically, as
detailed in Section 8. The right side is the indicator of a disjoint
union by (3); every term on the left is nonnegative. Thus the identity
also proves that (6) is a disjoint 27-cell set and that \(B_0+S\) is a
positive packing of exactly its stated complement. Adding
\(A\setminus A_0\) gives the actual initial scaffold on the full common
support.

\section{5. The eight-bone bridge and the completed endpoint
identity}\label{the-eight-bone-bridge-and-the-completed-endpoint-identity}

The parent first-residue edit sequence is an explicit signed bone vector
\(E_9\), with \[
 \partial E_9=\mathbf1_{Z_9}.
\] Its last upper patch, call it \(U\), has boundary \[
 \partial U=\mathbf1_{U_9}+\mathbf1_{\beta(c-1,b')},
 \quad c=m+2k-4,\quad b'=2k-5,
\] where \[
 \beta(x,y)=\{(x,y),(x+1,y),(x+1,y+1)\}.
\] This equation is the original eighteen-cell terminal partition in
\texttt{../turn9/PROOF.md}, not an abstract endpoint assertion.

The new bridge has anchor \[
 x=m+2k-4,\qquad y=2k-7.
\] Its old and new bones, in coordinates relative to \((x,y)\), are:

\begin{longtable}[]{@{}ll@{}}
\toprule\noalign{}
Old bones & New bones \\
\midrule\noalign{}
\endhead
\bottomrule\noalign{}
\endlastfoot
\(H(0,4),H(0,5),H(0,6)\) & \(V(0,0),V(0,3),V(1,0)\) \\
\(H(1,2),H(1,3),H(1,7)\) & \(V(1,3),V(2,2),V(2,5)\) \\
\(V(0,1),V(3,4)\) & \(V(3,2),V(3,5)\) \\
\end{longtable}

Let \(B_6\) denote new minus old. Expanding the original offsets shows
that both the old bones plus \(\beta(x,y)\), and the new bones plus
\(\beta(x,y+6)\), partition the same 27 cells. Consequently \[
 \partial B_6=\mathbf1_{\beta(x,y)}-
                     \mathbf1_{\beta(x,y+6)}.             \tag{8}
\] The relative coordinate map is translation by \((x,y)\), with
translation by \((-x,-y)\) as inverse. The two positive partitions give
inverse local replacement maps.

The original endpoint coordinates satisfy \[
 T_1\beta(c-1,b')=\beta(x,y),\qquad
 T_2 A_9=\beta(x,y+6).                                    \tag{9}
\] Define the complete second-residue edit chain \[
 E_2=T_1(E_9-U)+B_6+T_2E_9.                               \tag{10}
\] Substitute the three original boundaries and use (9): \[
\begin{split}
\partial E_2
={}&\mathbf1_{T_1(A_9\sqcup F_9\sqcup R_9)}
 -\mathbf1_{T_1\beta(c-1,b')}\\
 &+\mathbf1_{\beta(x,y)}-\mathbf1_{T_2A_9}
 +\mathbf1_{T_2Z_9}\\
={}&\mathbf1_{Z_2}.
\end{split}                                               \tag{11}
\] This is the endpoint cancellation that the interrupted attempt
lacked. Its two geometric changes are the eight-bone bridge and the
second deleted stone in (5).

The formula is also expanded independently by
\texttt{symbolic\_certificate.py}. It releases \(4m+12k+1\) old scaffold
bones and inserts \(4m+12k+10\) new bones, a gain of nine bones, exactly
the 27 filled cells. These are counts of the displayed repair, not an
optimality assertion.

\section{6. Original source ownership and
positivity}\label{original-source-ownership-and-positivity}

A signed identity does not supply the old generators needed to execute a
positive edit. For every negative generator in (10), \texttt{edits.json}
records an original owner of one of two types:

\begin{enumerate}
\def\labelenumi{\arabic{enumi}.}
\tightlist
\item
  a common source label \((p,y,\ell)\) from (3), proved to lie outside
  \(A_0\);
\item
  an exact family and index in the original parent \(B_8(m,k)\),
  followed by its actual rotation and the initial translation
  \((-1,-1)\).
\end{enumerate}

These are literal inverse index formulas. \texttt{ownership.py}
substitutes them into all three original cell coordinates and verifies
equality as an identity in \((m,k,i)\). Bounds are proved on the full
parameter domain. Every distinct pair of atoms that could share an owner
is checked against the two original label coordinates; their equality
produces a rational contradiction. Within a family, a nonzero index
coefficient proves injectivity.

The verifier reconstructs 1,479 exact affine certificates for the
uniform second residue, including 71 original generator identities and
733 potential duplicate-source comparisons. It uses only Python integer
and rational arithmetic. A certificate for a bound is a nonnegative
rational combination of the domain inequalities equal to that bound. A
contradiction certificate combines the domain and equality constraints
into \(-1\ge0\). No floating-point feasibility status is accepted as a
certificate.

The initial scaffold is a matching by (7), and the source certificate
proves that all negative edit generators are distinct members of that
initial scaffold. Subtracting them therefore leaves nonnegative
coefficients. Adding the positive edit generators leaves a nonnegative
tile vector. Equations (7) and (11) give its original boundary: \[
 \partial(B'+S)=\mathbf1_{\Omega_\Delta}+\partial A.        \tag{12}
\] The target coefficients are all one, so every tile is contained, all
tiles are disjoint, and every target cell is covered. This proves the
positive core tiling throughout the unbounded domain. No step assumes a
positive lifting theorem.

The source-to-tail map (4) now returns the core to every original
exterior benzel. Adjoining the frozen original bones gives \[
 (\Gamma_\Delta(d,h)\setminus A)\cup B'\cup S,              \tag{13}
\] and (1), (12) prove its positive partition. The counts are
\(t_m-3k-2\) stones and \(3h(h+d)\) bones.

\section{7. The four exceptional deletion
counts}\label{the-four-exceptional-deletion-counts}

The parent terminal prescription uses \(k-4\) vertical vacancy moves and
must not be used as a negative-length formula at \(k=2,3\). The
following exact constant patches replace it.

Let \(P\) be the parent edit sequence stopped immediately before its
\texttt{lower} group. Set \[
 c=m+2k-4,\qquad b=1-k.
\] For \(k=2,3\), the remaining twelve cells, relative to \((c,b)\), are
\[
\begin{split}
 H_k={}&\alpha(0,-3)\sqcup R(-3,0)\sqcup V(-1,3k-5)\\
 &\sqcup\{(2,3k-2),(2,3k-1),(3,3k-1)\}.
\end{split}\tag{14}
\] The fixed terminal vector \(Q_k\), printed in Appendix A, satisfies
\[
 \partial Q_k=\mathbf1_{H_k}.
\] For \(k=2\), it releases 11 bones and inserts 15. For \(k=3\), it
releases 14 and inserts 18. Every equality is proved by expanding the
literal finite offsets. Thus \(P+Q_k\) completes the first residue for
\[
 (q,m)=(7,m\ge9),\qquad(10,m\ge12).
\]

For the second residue, apply \(T_1P\), then a fixed direct bridge
\(D_k\), then \(T_2(P+Q_k)\). The direct bridge is anchored at \[
 (c',b')=(m+2k-3,-k-1).
\] In its relative coordinates it replaces \[
 \alpha(0,-3)\sqcup R(-3,0)
 \quad\text{by}\quad\beta(-1,3k).
\] For \(k=2\), it releases 16 bones and inserts 17; for \(k=3\), it
releases 20 and inserts 21. The complete tables are in Appendix A. The
resulting vacancy is exactly \(T_2A_9\), because its anchor is \[
 (m+2k-4,2k-1).
\] The endpoint cancellation is therefore the same original cell
cancellation as (11), with the direct bridge taking the place of the
invalid short lower/beta prescription.

This proves the second residue for \[
 (q,m)=(8,m\ge10),\qquad(11,m\ge13).
\] For these four branches, \(k\) is fixed but \(m\) is not bounded
above. The exact symbolic verifier specializes only \(k\), leaving the
original exterior parameter formal. Source ownership is proved by the
same affine mechanism on the domains \(k=k_0\), \(m\ge3k_0+r+2\). It
does not claim an identity for other \(k\)-values.

The four source verifications have, respectively, 407, 968, 559, and
1,178 exact certificates. The initial sources are retained and every old
edit generator remains distinct. The positivity proof of Section 6
applies to these literal tables. The smallest finite tail is included by
(4).

\section{8. Exact algebraic certificates and index
domains}\label{exact-algebraic-certificates-and-index-domains}

The cell-module map \[
 \mathbb Z[\mathbb Z^2]\longrightarrow
 \mathbb Z[X^{\pm1},Y^{\pm1}],\qquad[(x,y)]\mapsto X^xY^y
\] has inverse coefficient extraction. The four tile boundaries are
exactly \[
 1+X+Y,\quad1+X+X^2,\quad1+Y+Y^2,
 \quad1+XY^{-1}+X^2Y^{-2}.
\] Retain both independent parameters in \[
 \mathcal L=\mathbb Z[X^{\pm1},Y^{\pm1},
 M_x^{\pm1},M_y^{\pm1},K_x^{\pm1},K_y^{\pm1}].
\] The original specialization is \[
 M_x\mapsto X^m,\ M_y\mapsto Y^m,\
 K_x\mapsto X^k,\ K_y\mapsto Y^k.
\] Its kernel is generated by the four displayed parameter relations,
and the quotient inverse is the inclusion of \(X,Y\). This identifies
the exact relationship between formal parameter calculations and
original cells.

For an indexed anchor whose coordinates are affine in \(m,k,i\), the
finite sum is represented by its actual geometric progression. The
checker expands constant families directly. All denominators are nonzero
Laurent polynomials in \(X,Y\) alone; specialization never changes them.
The original Laurent rings are integral domains, so localization at
their product is injective.

For the new uniform branch, the scaffold numerator has 273 rational
input terms, 32,274 expanded numerator terms, and no surviving
coefficient. The endpoint identity has 228 rational input terms, 10,512
expanded terms, and no surviving coefficient. The separate fixed-\(k\)
expansions for the four exceptions also have zero numerator. These are
exact integer polynomial identities before choosing \(m\); the uniform
branch also retains \(k\).

Positive source domains are checked separately.
\texttt{source\_domains.py} covers both uniform scaffold tables, the
four fixed specializations, and the four retained fixed-deletion tables
\(q=1,2,4,5\). It verifies 1,291 exact affine bounds/contradictions:
nonnegative family lengths, original row/depth bounds, each prefix/tail
branch, and source distinctness. The four fixed-deletion polynomial
identities are replayed as well. Their construction data are preserved
in \texttt{../turn7}.

This makes the proof's logical order explicit: the polynomial identity
establishes original cell coefficients; the index certificates establish
that the expressions enumerate actual nonnegative generators; the source
owners establish that subtraction removes actual present tiles. Neither
a rational-function identity alone nor a finite tiling sample is
substituted for those three calculations.

\Needspace{22\baselineskip}

\section{9. Exhaustion of all canonical
parameters}\label{exhaustion-of-all-canonical-parameters}

Every \(q\in[0,m-2]\) has exactly one route in the following table.

\begin{longtable}[]{@{}
  >{\raggedright\arraybackslash}p{(\columnwidth - 4\tabcolsep) * \real{0.3333}}
  >{\raggedright\arraybackslash}p{(\columnwidth - 4\tabcolsep) * \real{0.3333}}
  >{\raggedright\arraybackslash}p{(\columnwidth - 4\tabcolsep) * \real{0.3333}}@{}}
\toprule\noalign{}
\begin{minipage}[b]{\linewidth}\raggedright
Canonical deletion count
\end{minipage} & \begin{minipage}[b]{\linewidth}\raggedright
Complete construction
\end{minipage} & \begin{minipage}[b]{\linewidth}\raggedright
Original domain
\end{minipage} \\
\midrule\noalign{}
\endhead
\bottomrule\noalign{}
\endlastfoot
\(q=0\) & Stable bands plus the all-stone core (2) & \(m\ge1\) \\
\(q=1,2,4,5\) & Retained literal fixed-deletion tables & \(m\ge q+2\) \\
\(q=3k,\ k\ge1\) & Retained threefold construction & \(m\ge3k+2\) \\
\(q=3k+1,\ k\ge4\) & Retained asymmetric construction & \(m\ge3k+3\) \\
\(q=3k+2,\ k\ge4\) & Sections 4--6 & \(m\ge3k+4\) \\
\(q=7,10\) & Fixed \(k=2,3\), first-residue terminal & \(m\ge q+2\) \\
\(q=8,11\) & Fixed \(k=2,3\), direct second-residue bridge &
\(m\ge q+2\) \\
\end{longtable}

To verify exhaustiveness, divide \(q\) by three. For remainder zero,
\(q=0\) or the threefold row applies. For remainder one, quotient
\(k=0,1,2,3\) gives \(q=1,4,7,10\), and all larger quotients use the
uniform first residue. For remainder two, quotients \(0,1,2,3\) give
\(2,5,8,11\), and all larger quotients use the new uniform second
residue. Every required lower bound on \(m\) is exactly the already
proved canonical inequality \(q\le m-2\).

Thus no canonical value is left to a search, unproved lifting statement,
or future endpoint. The original receiving parameter is \(d\ge m\), so
(13), or the corresponding retained literal construction, applies.

The two small \(d=2\) cases are also explicit. For \(h=0\), \(V(2,4)\)
is the central right stone. For \(h=1\), the same original three-sector
band formula gives the following nine bones:

\[
\begin{gathered}
D(-2,1),\quad D(-2,2),\quad D(-1,-1)\\
H(-3,3),\quad H(-1,2),\quad H(0,1)\\
V(1,-2),\quad V(2,-2),\quad V(3,-1)
\end{gathered}
\]

Their cell sets partition the original \(V(5,7)\), as direct offset
expansion shows. This completes the theorem in Section 1 for every
\(d\ge2\).

\section{10. Original Split-Zero complexes and their comparison
kernels}\label{original-split-zero-complexes-and-their-comparison-kernels}

This section records the concepts used in the construction, not a claim
that adjoining a zero alone supplies positive tilings. The source is the
retained Split-Zero reconstruction and internal-homology development at
the exact mathematical pin listed in \texttt{PROVENANCE.md}.

Fix one of the newly completed repair configurations. Let \(A^-\) be its
matching of old edit bones and \(A^+\) its matching of new edit bones,
and let \(Z\) be its actual initial vacancy. The proved endpoint
identity is \[
 \partial(A^+-A^-)=\mathbf1_Z.                            \tag{15}
\] The supports \(Z\) and the old cells are disjoint. Take \[
 L=\{\bot\}\sqcup\mathcal P(A^-),\qquad
 \Lambda_E=Z\sqcup\operatorname{cells}(E).
\] Bottom has empty physical support, whereas the nonbottom label
\(E=\varnothing\) has physical support \(Z\). The join map to cell
supports sends unions to unions and bottom to the empty set. Its inverse
on the image identifies which disjoint old bone cells are included.

At each supported label use the actual incidence complex of contained
bone generators \[
 C_E^0\xrightarrow{\partial_E}C_E^1\longrightarrow0,
 \qquad C_E^1=\mathbb Z[\Lambda_E].
\] The original inclusions intertwine differentials. Reconstruct the
totals \(\mathcal T(C^i)=\coprod_E C_E^i\) over
\(G(\mathbb Z)=\mathbb Z\sqcup\{\tau\}\). Their addition transports to
the union support. The supported scalar zero \(e=0_{\mathbb Z}\) and
absent zero act by \[
 e(E,x)=(E,0),\qquad \tau(E,x)=(\bot,0).
\] The differential square is the correctly typed map
\((E,n)\mapsto(E,0)\). Projection to the support proves its relationship
to the constant absent map; no label disappears when a boundary becomes
zero.

The internal homology projection \[
 (E,p)\longmapsto(E,[p])
\] coequalizes the actual incoming boundary and its supported-zero
companion. Its universal property follows directly: a split-linear map
equalizing those arrows takes two representatives differing by an
original boundary to the same value, using fibre-zero absorption. It
therefore descends uniquely; lifting representatives verifies addition,
scalar action, and naturality. No cancellation in an arbitrary target
semimodule is used.

\subsection{10.1 The augmented positive
fibre}\label{the-augmented-positive-fibre}

Let \(v_E=\mathbf1_{\Lambda_E}\) and \[
 \widehat C_E^0=\mathbb Z\oplus C_E^0,
 \qquad\widehat d_E(r,n)=\partial n-rv_E.
\] For \(E\subseteq E'\), let \(\eta_{EE'}\) be the sum of newly
released original bones. The original partition gives \[
 v_{E'}=jv_E+\partial\eta_{EE'}.
\] The genuine linear transition \[
 \widehat j_{EE'}(r,n)=(r,jn+r\eta_{EE'})
\] therefore satisfies
\(\widehat d_{E'}\widehat j_{EE'}=j\widehat d_E\). The identity
\(\eta_{EE''}=j\eta_{EE'}+\eta_{E'E''}\) proves composition.

Charge-one nonnegative integral cycles are exactly the original positive
patch tilings: \(\partial n=v_E\) forces coefficient one at each cell.
The inclusion into the full integral cycle module retains the same
charge and every original tile coefficient. Equation (15) supplies the
actual positive cycle \((1,A^+)\) at full support. The elementary
positive partitions and their inverses supply the concrete local maps
realizing that cycle.

\subsection{10.2 The full integral matching-comparison
kernel}\label{the-full-integral-matching-comparison-kernel}

At the fixed full support, write \(C=\mathbb Z[\Lambda]\),
\(B_-=\operatorname{im}\partial_{A^-}\), and
\(B_+=\operatorname{im}\partial_{A^+}\). The comparison is the original
representative map \[
 C/B_-\longrightarrow C/(B_-+B_+),
\] with kernel \[
 K\cong B_+/(B_-\cap B_+).                                \tag{16}
\] The map sends a new-boundary representative to its original class.
Its inverse writes a killed representative as \(a+b\), with
\(a\in B_-\), \(b\in B_+\), and sends it to \([b]\). Two decompositions
differ precisely by \(B_-\cap B_+\), proving well-definedness and both
inverse identities.

For completeness, the entire kernel is computed integrally. Build the
bipartite graph of actual old and new bones, with each shared original
cell as an edge. A singly owned cell marks its owner. Equality of old
and new boundary combinations forces coefficients equal along edges and
zero at marked vertices. They vanish on open components and have one
common integer value on each closed component. Conversely, a closed
component gives the equality of the sums of its original old and new
boundaries.

Those intersection generators are primitive and have disjoint supports.
Pick one new-bone pivot in each closed component. The remaining
new-boundary classes form an integral basis of (16). The coordinate map
subtracts the pivot coefficient on each closed component; the inverse
inserts zero at pivots. Original-cell forest paths give integer duals:
root each open component at ground and each closed component at its
pivot, and assign oriented path coefficients \(0,1,-1\) to its cell
edges. Summing at old vertices gives zero, and summing at new vertices
gives the Kronecker evaluations of the displayed basis.

Any chosen vacancy cell evaluates \([\mathbf1_Z]\) to one and
annihilates every old boundary. It supplies a split copy of
\(\mathbb Z\) before the comparison. Equation (15) kills that
distinguished class at its own support. This one summand is retained
alongside the full kernel, not substituted for it.

The selected matching complexes also map into the complex of all
contained bones by literal inclusion on tile generators and identity on
cells. The further quotient kernel is \[
 \operatorname{im}\partial_{\rm all}/(B_-+B_+).
\] This records precisely the relationship with the unrestricted
incidence complex; no vanishing of that additional kernel is assumed.
None of these elementary matching calculations is claimed to be a new
computation of the known global bone homology group.

\section{11. Returning to the original Problem
7}\label{returning-to-the-original-problem-7}

The retained source-matching note \texttt{reference-earlier-proofs.md},
Section 3, specifies the original admissible domain \(a,b\ge2\),
\(a\le2b\), \(b\le2a\), and the three residue cases of the
Conway--Lagarias invariant. That note is included unchanged and its
historical status statements remain historical. Reflection permits
\(a\le b\) through the exact involution in Section 2.

For \(a+b\equiv0\pmod3\), the inverse parameters are \(d=b-a\),
\(h=(2a-b)/3\), and \(\Delta=t_d-h\). Admissibility gives \(h\ge0\), and
nonnegative invariant gives \(h\le t_d\). The possibilities \(d=0,1\)
would force \(h=0\) and violate \(a\ge2\). Thus \(d\ge2\), and the
completed constructor applies.

For \(a+b\equiv1\pmod3\), there is a direct all-right-stone
construction. Put \(A=1-a,B=b-1,r=2-a\pmod3\). For an original
barycentric cell \(v\), consider the three sum-zero anchors \(u=v-e_j\)
of containing right stones. Their first differences occupy all three
residues; choose the unique anchor in residue \(r\). All differences of
a sum-zero anchor are congruent modulo three. They lie initially in
\([A-1,B+1]\), but the endpoints \(A+1,B-1\) both have residue \(r\),
and none of \(A-1,A,B,B+1\) does. Hence the chosen anchor differences
lie in \([A+1,B-1]\), and its three cells lie in the original benzel.
Every cell of that stone selects the same anchor. The map from selected
anchors and offset indices to original cells is therefore bijective,
with the displayed selection as inverse. \texttt{residue\_one}
implements this construction directly.

For \(a+b\equiv2\pmod3\), the retained source-matching note invokes the
already published generalized-compression Theorem 1.1 at
arXiv:2403.07663v1, in the exact parameters \[
 n=b+1-a,\qquad j=(2a-b-1)/3,
 \qquad(a,b)=(n+3j,2n+3j-1).
\] The inverse formulas are displayed and \(n\ge1,j\ge0\) follow from
admissibility and the residue. That established result gives positive
tiling counts using the original right stone and two original bone
orientations. The literal inclusion of its tile set into the type-103
set preserves each placement and its three cells; its induced incidence
square commutes and its map on tilings is injective. This is the exact
published input of the full-P7 assembly. The present package does not
assert a self-contained new construction or independent proof audit of
that compression result.

Accordingly, the previously recorded route to the full existence
statement now has its remaining residue-zero case supplied by a complete
original-coordinate construction. This is independent of Sharma's
reported manuscript in method and proof dependencies, while making no
claim of first proof or novel bare existence.

\section{12. Verification, limitations, and
reproduction}\label{verification-limitations-and-reproduction}

The mathematical proof uses parameter-preserving polynomial identities,
exact affine certificates, the retained full packing proofs, and the
positive incidence arguments above. Finite tests are additional checks
of the implementation, not the source of an unbounded quantifier.

Run from this directory:

\begin{Shaded}
\begin{Highlighting}[]
\ExtensionTok{python}\NormalTok{ verify.py }\AttributeTok{{-}{-}max{-}d}\NormalTok{ 12 }\AttributeTok{{-}{-}max{-}k}\NormalTok{ 9 }\AttributeTok{{-}{-}output}\NormalTok{ evidence/normal.json}
\ExtensionTok{python} \AttributeTok{{-}O}\NormalTok{ verify.py }\AttributeTok{{-}{-}max{-}d}\NormalTok{ 12 }\AttributeTok{{-}{-}max{-}k}\NormalTok{ 9 }\AttributeTok{{-}{-}output}\NormalTok{ evidence/optimized.json}
\end{Highlighting}
\end{Shaded}

The verifier uses only the Python standard library. Optional scripts
that discovered rational multipliers or fixed finite patches are not
invoked. The source map, all small constant tables, old-generator
inverse maps, exact certificates, and retained mathematical dependencies
are included.

A complete run checks every allowed \(h\) at \(d=2,\ldots,12\), its
reflection, larger uniform second-residue cores at their smallest and
larger \(m\), all four exceptional families at three \(m\)-values and
two receiving values, all duals in eight original matching comparisons,
actual augmented transitions, and deliberately corrupted inputs. Exact
counts and file identities belong to the emitted receipts.

There is no first-discovery claim, no new Lean build certificate, and no
independent specialist acceptance. This manuscript closes the displayed
construction and its canonical parameter coverage. The separate prior
literature claims and the published residue-two input keep their own
provenance and verification scopes.

\section{Appendix A. The four constant terminal
tables}\label{appendix-a.-the-four-constant-terminal-tables}

All coordinates in these tables are relative to the anchors in Section
7. The table entries are original placed bones, not quotient
representatives. Old/new equalities follow by literal offset expansion;
\texttt{small\_symbolic.py} verifies the aggregate unbounded-\(m\)
identities and the ownership certificates supply the actual source
indices.

\Needspace{24\baselineskip}

\subsection{\texorpdfstring{Terminal
\(Q_2\)}{Terminal Q\_2}}\label{terminal-q_2}

Original releases: 11 bones. Original insertions: 15 bones.

\begin{longtable}[]{@{}ll@{}}
\toprule\noalign{}
Old placed bone & New placed bone \\
\midrule\noalign{}
\endhead
\bottomrule\noalign{}
\endlastfoot
\(H(-1,0)\) & \(D(0,-3)\) \\
\(H(0,1)\) & \(D(0,-2)\) \\
\(H(0,2)\) & \(D(0,0)\) \\
\(H(0,3)\) & \(D(0,1)\) \\
\(V(-2,1)\) & \(D(1,-2)\) \\
\(V(1,-3)\) & \(H(-3,0)\) \\
\(V(2,-5)\) & \(H(-3,1)\) \\
\(V(2,-2)\) & \(H(-2,2)\) \\
\(V(3,-4)\) & \(H(-2,3)\) \\
\(V(3,-1)\) & \(V(1,1)\) \\
\(V(3,2)\) & \(V(2,0)\) \\
& \(V(2,3)\) \\
& \(V(3,-3)\) \\
& \(V(3,0)\) \\
& \(V(3,3)\) \\
\end{longtable}

\Needspace{26\baselineskip}

\subsection{\texorpdfstring{Direct bridge
\(D_2\)}{Direct bridge D\_2}}\label{direct-bridge-d_2}

Original releases: 16 bones. Original insertions: 17 bones.

\begin{longtable}[]{@{}ll@{}}
\toprule\noalign{}
Old placed bone & New placed bone \\
\midrule\noalign{}
\endhead
\bottomrule\noalign{}
\endlastfoot
\(H(-1,0)\) & \(D(0,-3)\) \\
\(H(-1,6)\) & \(D(0,-2)\) \\
\(H(0,1)\) & \(D(0,0)\) \\
\(H(0,2)\) & \(D(0,1)\) \\
\(H(0,3)\) & \(D(1,-2)\) \\
\(H(0,7)\) & \(H(-3,0)\) \\
\(V(-2,1)\) & \(H(-3,1)\) \\
\(V(-1,1)\) & \(H(-2,2)\) \\
\(V(1,-3)\) & \(H(-2,3)\) \\
\(V(2,-5)\) & \(H(1,6)\) \\
\(V(2,-2)\) & \(H(1,7)\) \\
\(V(2,4)\) & \(V(1,1)\) \\
\(V(3,-4)\) & \(V(2,0)\) \\
\(V(3,-1)\) & \(V(2,3)\) \\
\(V(3,2)\) & \(V(3,-3)\) \\
\(V(3,5)\) & \(V(3,0)\) \\
& \(V(3,3)\) \\
\end{longtable}

\Needspace{27\baselineskip}

\subsection{\texorpdfstring{Terminal
\(Q_3\)}{Terminal Q\_3}}\label{terminal-q_3}

Original releases: 14 bones. Original insertions: 18 bones.

\begin{longtable}[]{@{}ll@{}}
\toprule\noalign{}
Old placed bone & New placed bone \\
\midrule\noalign{}
\endhead
\bottomrule\noalign{}
\endlastfoot
\(H(-1,0)\) & \(D(0,-3)\) \\
\(H(1,4)\) & \(D(0,-2)\) \\
\(V(-2,1)\) & \(D(1,-2)\) \\
\(V(-1,1)\) & \(D(1,-1)\) \\
\(V(0,2)\) & \(D(2,-1)\) \\
\(V(1,-3)\) & \(H(-3,0)\) \\
\(V(2,-5)\) & \(H(-3,1)\) \\
\(V(2,-2)\) & \(H(-2,2)\) \\
\(V(3,-4)\) & \(H(-2,3)\) \\
\(V(3,5)\) & \(H(0,0)\) \\
\(V(4,-3)\) & \(H(0,4)\) \\
\(V(4,0)\) & \(H(2,7)\) \\
\(V(4,3)\) & \(H(2,8)\) \\
\(V(4,6)\) & \(V(-1,4)\) \\
& \(V(3,4)\) \\
& \(V(4,-2)\) \\
& \(V(4,1)\) \\
& \(V(4,4)\) \\
\end{longtable}

\Needspace{30\baselineskip}

\subsection{\texorpdfstring{Direct bridge
\(D_3\)}{Direct bridge D\_3}}\label{direct-bridge-d_3}

Original releases: 20 bones. Original insertions: 21 bones.

\begin{longtable}[]{@{}ll@{}}
\toprule\noalign{}
Old placed bone & New placed bone \\
\midrule\noalign{}
\endhead
\bottomrule\noalign{}
\endlastfoot
\(H(-1,0)\) & \(D(0,-3)\) \\
\(H(-1,9)\) & \(D(0,-2)\) \\
\(H(0,1)\) & \(D(0,0)\) \\
\(H(0,5)\) & \(D(0,1)\) \\
\(H(0,6)\) & \(D(1,-2)\) \\
\(H(0,10)\) & \(H(-3,0)\) \\
\(H(1,2)\) & \(H(-3,1)\) \\
\(H(1,3)\) & \(H(-2,2)\) \\
\(H(1,4)\) & \(H(-2,3)\) \\
\(V(-2,1)\) & \(H(1,9)\) \\
\(V(-1,1)\) & \(H(1,10)\) \\
\(V(0,2)\) & \(V(0,4)\) \\
\(V(1,-3)\) & \(V(1,1)\) \\
\(V(2,-5)\) & \(V(1,4)\) \\
\(V(2,-2)\) & \(V(2,0)\) \\
\(V(2,7)\) & \(V(2,3)\) \\
\(V(3,-4)\) & \(V(2,6)\) \\
\(V(3,-1)\) & \(V(3,-3)\) \\
\(V(3,5)\) & \(V(3,0)\) \\
\(V(3,8)\) & \(V(3,3)\) \\
& \(V(3,6)\) \\
\end{longtable}

\section{Appendix B. Literal indexed edit
tables}\label{appendix-b.-literal-indexed-edit-tables}

\texttt{edits.json} contains the complete uniform second-residue table.
The four exceptional tables have filenames \texttt{edits-kK-rR.json},
with \((K,R)=(2,1),(2,2),(3,1),(3,2)\). They retain all their original
owner maps after fixing \(k\). No unspecified geometric pattern is
called by the constructor.

A row stores a kind and two vectors \([a,b,c,e]\); its anchor coordinate
is exactly \(am+bk+ci+e\). A count vector \([a,b,e]\) means
\(0\le i<am+bk+e\). Common owners store their original sector, row,
depth, and triangular branch. Parent owners store the exact part,
family, inverse indices, and rotation. Their equality is verified on all
three original cell coordinates.

The retained full first-residue edit table gives the literal meaning of
\(E_9\) and \(U\) in (10). \texttt{build\_edits.py} records their exact
affine transport; the final constructor reads the fully expanded table
and calls no discovery script.

\end{document}
